\section{The Blessed and the Justice of God}

The question of those who may finally refuse God touches the hope of heaven with particular pain. It becomes most difficult when we think of someone we love. The blessed possess perfect happiness and are wholly united to God's will; eternal loss remains a real possibility for the creature who dies in unrepented mortal sin. We cannot resolve this gravity by imagining that charity becomes indifference, or that God's justice becomes an evil from which the blessed must be protected. His goodness is the ground of their happiness, and their love is perfected in agreement with him.

Aquinas addresses the question in the \emph{Supplement} to the \emph{Summa} (q.~94, a.~3). He says that the blessed rejoice in divine justice and in their own deliverance. He distinguishes this from rejoicing directly in another's suffering as suffering, which would belong to hatred. The punishment enters his account indirectly, under the good of justice; cruelty is excluded. His explanation is demanding, and its meaning is lost if it is made either a permission for malice or a promise that judgment will cease to matter. The blessed love the holy order of God with a love entirely free from sin.

This account does not give us a description of the precise experience of a mother, husband, or friend faced with another person's loss. Our present emotional imagination cannot supply what has not been revealed. We know that the blessed will be neither deceived nor miserable, and that their charity will contain no hatred. We also know that we are not presently entrusted with declaring the final destiny of particular people. The proper response to someone whose life troubles us is prayer, truthful witness, patient love, and such help as charity allows. The time given to us is a time in which we seek salvation together.

Concern for another's salvation can become anxious when we assume responsibility for choices we cannot make for him. His life matters immeasurably, but we cannot believe, repent, or persevere in his place. We can be faithful to our own responsibility and entrust him to a mercy wiser than our fear. That entrustment does not end intercession. It purifies it, freeing prayer from the demand that God disclose beforehand what only he can judge. We ask for his salvation because we desire his good, and we continue to love him while leaving his freedom and final judgment in God's hands.

The consideration of divine justice should make our desire for heaven more charitable now. A person who enjoys imagining another's condemnation is consenting to something from which heaven would have to heal him. The saints' joy is the joy of God's goodness, and our preparation must move toward that goodness. We can oppose evil, seek protection for those harmed, and still pray for the conversion of the wrongdoer. Seeking protection and seeking conversion can place difficult demands upon us at the same time. We need God's help to fulfill both without excusing the wrong or wishing evil to its author.

